\documentclass[10pt,a4paper]{amsart}
\usepackage[margin=19mm]{geometry}
\usepackage{amsmath,amssymb,mathtools}
\usepackage[scr=rsfs,cal=euler]{mathalfa}
\usepackage{xfrac}
\usepackage[unicode,hidelinks,bookmarks=false]{hyperref}
\newcommand{\ie}{i.e.\ }
\newcommand{\cf}{cf.\ }
\newcommand{\resp}{respectively\ }
\newcommand{\Subsection}{\subsection}
\providecommand{\nobreakdash}{\nobreak\hbox{-}}
\setlength{\emergencystretch}{2em}
\multlinegap0pt
\usepackage{bm}
\usepackage{bbm}
\usepackage{dsfont}
\usepackage{xspace}
\usepackage{cancel}
\usepackage{mathtools}
\usepackage[shortlabels]{enumitem}
\usepackage{amsthm}
\makeatletter
\@ifundefined{thm}{\newtheorem{thm}{Thm}}{}
\@ifundefined{defi}{\newtheorem{defi}[thm]{Definition}}{}
\@ifundefined{lemma}{\newtheorem{lemma}[thm]{Lemma}}{}
\@ifundefined{prop}{\newtheorem{prop}[thm]{Proposition}}{}
\makeatother
\makeatletter
\def\Chat{\widehat{\mathbb{C}}}
\def\DD{\mathbb{D}}
\expandafter\ifx\csname named\endcsname\relax
\newenvironment{named}[1]
{\def\namedthmname{#1}%
	\refstepcounter{namedthm}%
	\namedthm\def\@currentlabel{#1}}
{\endnamedthm}
\fi
\renewcommand{\tilde}{\widetilde}
\makeatother
\title[Boundary behaviour of universal covering maps]{Boundary behaviour of universal covering maps}
\author{Gustavo R. Ferreira, Anna Jov\'e}
\date{Source: arXiv:2409.01070v1 (CC BY 4.0)}
\begin{document}
\maketitle

\noindent\emph{Source and licence.} Boundary behaviour of universal covering maps. Version arXiv:2409.01070v1 (CC BY 4.0), distributed under the Creative Commons Attribution 4.0 International licence, as recorded in the arXiv licence metadata for this exact version and shown on the abstract page https://arxiv.org/abs/2409.01070v1. Full rights notice: licenses/arxiv-2409.01070-CC-BY-4.0.txt.\par\smallskip

\noindent\textit{Editorial scope.} Counted unit: the Prop whose source label is \texttt{prop-associated-boundary-comp} in the article (source0.tex lines 1091--1096, source label \texttt{prop-associated-boundary-comp}), delivered together with the article's own complete original proof of it (source0.tex lines 1097--1111, one \texttt{proof} environment of 15 lines). No mathematics is written, repaired or re-derived: statement and proof are the article's own text line for line. Required context retained uncounted: prop:hypdisc (lines 319--325); lem:restriction (lines 476--479); lem:superSP (lines 547--561); lemma-lifts-of-curves (lines 735--759); prop-escaping-bounded-bungee (lines 920--929); lemma-equivalence-ohtsuka (lines 1072--1074); def:associatedaddress (lines 1087--1089). Every definition, display and label of those lines is retained unchanged. Omitted from this package and not redistributed: 1--318, 326--475, 480--546, 562--734, 760--919, 930--1071, 1075--1086, 1090--1090, 1112--1742. Nothing inside a retained excerpt is omitted or altered: every retained body is delivered exactly as the article's lines are written, so no omission receipt of any kind accompanies this package. Citations: the retained excerpts carry no citation command at all, so no bibliography entry of the article is transcribed and no reference is redistributed. Reference adaptation: none was needed and none was made; every reference inside the delivered unit resolves inside it, so no unresolved reference or citation is produced. Printed slips kept literal: no printed word, symbol, hypothesis or number of the counted statement or of its proof was altered. Layout and class: the article's own class is \texttt{amsart} with packages including inputenc, amsmath, amssymb, amsthm, mathtools, enumerate, xfrac, footmisc, hyperref, enumitem, babel, graphicx, subcaption, xcolor; the wrapper is the lane's pinned \texttt{amsart} editorial wrapper at 10pt on A4, which restates the theorem-like environments the retained text needs and adds the article's own macro definitions that the retained text needs (\texttt{\textbackslash Chat}, \texttt{\textbackslash DD}, \texttt{\textbackslash tilde}), with extra wrapper packages (bm, bbm, dsfont, xspace, cancel, mathtools, enumitem). The delivered numbering is the unit's own. Assets: no retained span contains a figure environment, a graphics-inclusion command, a PDF-page inclusion, a table or a TikZ picture; no image file is copied into this package, the package declares no asset dependency and no image file is redistributed with it. Licence: version arXiv:2409.01070v1 of this article is distributed under the Creative Commons Attribution 4.0 International licence, as recorded in the arXiv licence metadata for this exact version and shown on the abstract page https://arxiv.org/abs/2409.01070v1. A full rights notice is delivered at licenses/arxiv-2409.01070-CC-BY-4.0.txt. Characters: every retained excerpt is pure ASCII, so the delivered engine, XeLaTeX with its default Latin Modern text font, reports no missing character.
\begin{prop}{\bf (Properties of the hyperbolic metric in $ \mathbb{D} $)}\label{prop:hypdisc} The following hold.
	\begin{enumerate}[label={\em (\alph*)}]
		\item For all  $\gamma \in \textrm{\em Aut}(\mathbb{D})  $ and $ z,w\in\mathbb{D} $, we have $ \textrm{\em dist}_\mathbb{D}(z,w)= \textrm{\em dist}_\mathbb{D}(\gamma(z),\gamma (w)) $.
		\item The hyperbolic distance is a complete distance inducing on $ \mathbb{D} $ the Euclidean topology. More precisely, hyperbolic disks are Euclidean disks with (possibly) different center and radius.
		\item The geodesics for the hyperbolic distance are the Euclidean diameters of $ \mathbb{D} $ and the intersections with $ \mathbb{D} $ of Euclidean circles orthogonal to $ \partial\mathbb{D} $. In particular, geodesics are smooth curves and any two distinct points of $ \mathbb{D} $ are connected by a unique geodesic.
	\end{enumerate}
\end{prop}

\begin{lemma}{\bf (Restrictions of the universal covering)}\label{lem:restriction}
	Let $ \Omega $ be a hyperbolic domain, and let $ \pi\colon\mathbb{D}\to \Omega $ be a universal covering. If $ V\subset \Omega $ and $ U $ is a connected component of $ \pi^{-1}(V) $ and $\pi_1(V)\subset \pi_1(\Omega)$, then $ U $ is simply connected and
	\[\pi|_U\colon U\to V\] is a universal covering of $ V $.
\end{lemma}

\begin{lemma}{\bf (Properties of the hyperbolic metric in $ \Omega $)}\label{lem:superSP}
Let $\Omega\subset\Chat$ be a hyperbolic domain, and let $\pi\colon\DD\to\Omega$ be a universal covering. Consider the hyperbolic metrics $ \rho_\mathbb{D} $ and $ \rho_\Omega $ in $ \mathbb{D} $ and $ \Omega $, respectively. The following hold. 
	\begin{enumerate}[label={\em (\alph*)}]
		\item  $ \pi$ is a local isometry between $ (\mathbb{D}, \rho_\mathbb{D}) $ and $ (\Omega, \rho_\Omega) $.
		\item Let $z, w\in\Omega$. Then, for any $ \tilde z \in\pi^{-1}(z)$, we have
		\[ \textrm{\em dist}_\Omega(z,w)=\min \left\lbrace \textrm{\em dist}_\DD(z, \gamma(w))\colon \ \gamma \in\Gamma \right\rbrace .\]
		In particular, there exists $\tilde w\in\pi^{-1}(w)$ such that
		\[ \textrm{\em dist}_\Omega(z, w) = \textrm{\em dist}_\DD(\tilde z, \tilde w). \]
		\item $(\Omega, \rho_\Omega)$ is a complete metric space.
		\item Any two points $z$ and $w$ in $\Omega$ can be joined by a distance-minimizing geodesic. Furthermore, any two geodesics joining $z$ and $w$ are not homotopic to each other.
		\item Let $\left\lbrace z_n\right\rbrace _n\subset\Omega$ be a sequence such that $z_n\to\partial\Omega$ as $n\to+\infty$. Then, for any $r > 0$,
		\[ \mathrm{diam}_{\widehat{\mathbb{C}}}(D_\Omega(z_n, r))\to 0 \text{ as $n\to+\infty$,} \]
		where $\mathrm{diam}_{\widehat{\mathbb{C}}}$ denotes the spherical diameter.
		\end{enumerate}
\end{lemma}

\begin{lemma}{\bf (Lifts of non-contractible curves)}\label{lemma-lifts-of-curves}
	Let $ \Omega \subset \widehat{\mathbb{C}}$ be a multiply connected domain of connectivity greater than two, and let $ \pi\colon\mathbb{D}\to \Omega $ be a universal covering. Let $ \eta $ be a non-contractible Jordan curve in $ \Omega $, and let $ D_\eta $ be a connected component of $ \Omega\smallsetminus\eta $. Exactly one of the following holds. \begin{enumerate}[label={\em (\alph*)}] 
			\item  $D_\eta $ is conformally isomorphic to $ \mathbb{D}^* $, and the following hold.
		\begin{enumerate}[label={\em a.\arabic*}]
			\item $ \pi^{-1}(\eta) $ consists of  countably many disjoint degenerate crosscuts in $ \mathbb{D} $, whose endpoints are parabolic fixed points. Any two crosscuts in $ \pi^{-1}(\eta) $ have no  common endpoint.
			\item $ \pi^{-1}(D_\eta) $ consists of countably many simply connected domains, each of which touches $\partial\DD$ at exactly one point. Each connected component of $ \pi^{-1}(D_\eta) $ is bounded by one crosscut in $ \pi^{-1}(\eta) $.
			\item There exist curves $\eta^*\subset\Omega$ homotopic to $\eta$, such that $ \pi^{-1}(D_{\eta^*}) $ are horodisks.
		\end{enumerate}
		\item  $ D_\eta $ is conformally isomorphic to a non-degenerate annulus, and the following hold.
		\begin{enumerate}[label={\em b.\arabic*}]
			\item There exists a unique closed geodesic $\eta^*\subset\Omega$ homotopic to $\eta$.
				\item $ \pi^{-1}(\eta) $ consists of  countably many disjoint non-degenerate crosscuts in $ \mathbb{D} $, whose endpoints are hyperbolic fixed points. Any two crosscuts in $ \pi^{-1}(\eta) $ have no endpoints in common.
			\item $ \pi^{-1}(D_\eta) $ consists of countably many simply connected domains, each of which is unbounded in $ \mathbb{D} $. Each connected component of $ \pi^{-1}(D_\eta) $ is bounded by one crosscut in $ \pi^{-1}(\eta) $, i.e. it is a crosscut neighbourhood for that crosscut.
		\end{enumerate}
		\item Otherwise, $ D_\eta $ has connectivity greater than two, and the following hold.
				\begin{enumerate}[label={\em c.\arabic*}]
			\item There exists a unique closed geodesic $\eta^*\subset\Omega$ homotopic to $\eta$.
			\item $ \pi^{-1}(\eta) $ consists of  countably many disjoint non-degenerate crosscuts in $ \mathbb{D} $, whose endpoints are hyperbolic fixed points. Any two crosscuts in $ \pi^{-1}(\eta) $ have no endpoints in common.
			\item $ \pi^{-1}(D_\eta) $ consists of countably many simply connected domains, each of which is unbounded in $ \mathbb{D} $. Each connected component of $ \pi^{-1}(D_\eta) $ is bounded by countably many crosscuts in $ \pi^{-1}(\eta) $.
			\item For any crosscut in $ \pi^{-1}(\eta) $, each of its crosscut neighbourhoods contain infinitely many crosscuts in $ \pi^{-1}(\eta) $.
		\end{enumerate}
	\end{enumerate}

\noindent Moreover, let $ \eta_0 $ be another non-contractible curve in $ \Omega $. If $ \eta_0 $ is homotopic to $ \eta $ in $ \Omega $,  then for each crosscut in $ \pi^{-1}(\eta) $, there is a crosscut in $ \pi^{-1}(\eta_0) $  with the same endpoints in $ \partial\mathbb{D} $, and vice versa. If $ \eta_0 $ is not homotopic to $ \eta $ in $ \Omega $, then any two crosscuts in $ \pi^{-1}(\eta) $ and in $ \pi^{-1}(\eta_0) $ have no endpoints in common.
\end{lemma}

\begin{prop}{\bf (Boundary behaviour of non-tangential curves)}\label{prop-escaping-bounded-bungee}	Let $ \Omega $ be a multiply connected domain,  let $ \pi\colon\mathbb{D}\to \Omega $ be the universal covering and let $ e^{i\theta}\in\partial\mathbb{D} $. 
	Let $ \eta $ be a curve landing non-tangentially at $ e^{i\theta}$.
		\begin{itemize}
		\item  If $ e^{i\theta}\in\partial \mathbb{D} $ is of {escaping} type, then $ \pi(\eta (t))\to\partial \Omega $, as $ t\to 1^{-} $, and
		\[ Cl_{\mathcal{A}}(\pi, e^{i\theta})= Cl_{R}(\pi, e^{i\theta})\subset \partial \Omega. \] In particular, $ Cl_{\mathcal{A}}(\pi, e^{i\theta}) $ is closed and contained in a unique boundary component of $\Omega$.
		\item If  $ e^{i\theta}\in\partial \mathbb{D} $ is of bounded type, then $\left\lbrace  \pi(\eta (t))\colon t\in \left[ 0,1\right) \right\rbrace  $ is compactly contained in $ \Omega $. Moreover,
		\[ Cl_{\mathcal{A}}(\pi, e^{i\theta})=  \Omega. \] In particular, $ Cl_{\mathcal{A}}(\pi, e^{i\theta}) $ is open.
		\item If $ e^{i\theta}\in\partial \mathbb{D} $ is of {bungee} type, then $\left\lbrace  \pi(\eta (t))\colon t\in \left[ 0,1\right) \right\rbrace  $ is neither compactly contained in $ \Omega $ nor converging to $ \partial \Omega $.
	\end{itemize}
\end{prop}

\begin{lemma}{\bf (Depth does not depend on the exhaustion)}\label{lemma-equivalence-ohtsuka}
	Let $ \Omega $ be a multiply connected domain, and let $ \pi\colon \mathbb{D}\to\Omega $ be a universal covering. Let $ e^{i\theta} \in\partial\mathbb{D}$. Then, $ e^{i\theta} $ is of finite, infinite or oscillating depth  not depending on the chosen Ohtsuka exhaustion of $ \Omega $.
\end{lemma}

\begin{defi}{\bf (Associated boundary address)}\label{def:associatedaddress}
Let $e^{i\theta}\in\partial\DD$, and let $\Sigma$ be a component of $\partial\Omega$. We say that $e^{i\theta}$ has $\underline{s}(\Sigma) = \left\lbrace s_n\right\rbrace _n$ as an \textit{associated boundary address} if we can find a subsequence of fundamental crosscuts $\left\lbrace C_{m_n}\right\rbrace _n$ at $e^{i\theta}$ such that $\pi(C_{m_n}) = \sigma_{s_1s_2\dots s_n}$, for all $n\geq 1$.
\end{defi}

\begin{prop}{\bf (Associated boundary components)}\label{prop-associated-boundary-comp}
Let $ \Sigma $ be a boundary component of $ \Omega $ with address $ \underline{s} (\Sigma)$.
Then,	$ e^{i\theta} \in\partial\mathbb{D}$ has $ \underline{s} (\Sigma)$ as an associated boundary address if and only if it has $\Sigma$ as an associated boundary component. 

\noindent Moreover, if $ e^{i\theta} $ is of escaping type, then $  Cl_{\mathcal{A}}(\pi, e^{i\theta})\subset\Sigma$ and $ \Sigma $ is the unique boundary component associated to $ e^{i\theta}$. For every countable collection of non-isolated boundary components $ \left\lbrace \Sigma_n \right\rbrace _n$, there exist uncountably many points $ e^{i\theta}\in\partial\mathbb{D} $ of bungee type such that each $e^{i\theta}$ has $ \left\lbrace \Sigma_n \right\rbrace _n$ as associated boundary components.
\end{prop}

\begin{proof}
	First, we argue that we can restrict our attention to the radial limit set. Indeed, let $\left\lbrace z_n\right\rbrace _n\subset\DD$ be a sequence tending non-tangentially to $e^{i\theta}$ and such that $\pi(z_n)\to\partial\Omega$ as $n\to+\infty$. Then, by the definition of a hyperbolic Stolz angle, there exists $r > 0$ and a sequence $\left\lbrace w_n\right\rbrace _n\subset R_\theta$ such that $\textrm{dist}_\Omega(\pi(w_n), \pi(z_n)) < r$ for all $n$. By Lemma \ref{lem:superSP}, the spherical distance $\textrm{dist}_{\Chat}(z_n, w_n)$ tends to $ 0$ as $n\to+\infty$, and therefore both sequences accumulate at the same points on $ \partial\Omega $.

	Now, consider an exhaustion $ \left\lbrace \Omega_n\right\rbrace _n $ of $ \Omega $ such that all the curves in $ \left\lbrace \partial\Omega_n\right\rbrace _n $ are geodesics (which we can do by Lemmas \ref{lemma-lifts-of-curves} and \ref{lemma-equivalence-ohtsuka}). Let $ \left\lbrace C_m\right\rbrace _m $ be the sequence of fundamental crosscuts at $ e^{i\theta} $ given by the exhaustion  $ \left\lbrace \Omega_n\right\rbrace _n $. Then, $ R_\theta $ intersects each crosscut $ C_m $ exactly once, and intersects no other crosscut obtained from the exhaustion, since by Proposition \ref{prop:hypdisc} each curve $\pi^{-1}(\sigma_{s_0s_1\dots s_n})$ is a circle arc orthogonal to $\partial\DD$. 
	
	First assume $ e^{i\theta} \in\partial\mathbb{D}$ has $ \underline{s} (\Sigma)$ as an associated boundary address. By Definition \ref{def:associatedaddress}, there exists a subsequence $ \left\lbrace C_{m_n}\right\rbrace _n $ such that $ \pi(C_{m_n})=\sigma_{s_1s_2\dots s_n} $. This already implies that \[ Cl_R(\pi, e^{i\theta})\cap\Sigma \neq\emptyset.\] For the other implication, we claim $ \pi (R_\theta ) $ cannot accumulate in a boundary component whose boundary address is not associated to $ e^{i\theta} $. Indeed, that would require $\pi(R_\theta)$ to cross a boundary curve $\sigma_{s_1s_2\dots s_n}'$ whose lifts are, by hypothesis, not part of the fundamental crosscuts at $e^{i\theta}$, and therefore cannot intersect $R_\theta$.
	
	The statement for points of escaping type follows straightforward from Proposition \ref{prop-escaping-bounded-bungee}. 

	It is left to see that for every countable collection of non-isolated boundary components $ \left\lbrace \Sigma_n \right\rbrace _n$, there exist uncountably many  points $ e^{i\theta}\in\partial\mathbb{D} $ of bungee type associated to $ \left\lbrace \Sigma_n \right\rbrace _n$. Equivalently, we shall prove that, given a sequence $ \left\lbrace \Sigma_n \right\rbrace _n$ of non-isolated boundary components, there exist uncountably many points $ e^{i\theta}\in\partial\mathbb{D} $ having $ \left\lbrace \underline{s}(\Sigma_n) \right\rbrace _n$ as associated boundary addresses. To do so, we rely on the nestedness properties of the crosscuts coming from the exhaustion. We will outline the argument for two non-isolated boundary components $\Sigma$ and $\Sigma'$; the argument for countably many components is similar (one can use a diagonal argument to alternate between components). Indeed, let $\underline{s}(\Sigma) = s_1s_2\dots$ and $\underline{s}'=\underline{s}(\Sigma') = s_1's_2'\dots$ denote the boundary addresses of $\Sigma$ and $\Sigma'$, where $\Sigma, \Sigma'\subset\partial\Omega$ are non-isolated boundary components, and let $D_{s_1s_2\dots s_n}$ and $D_{s_1's_2'\dots s_n'}$ denote the complementary components of $\Omega_n$ such that
	\[ \Sigma = \bigcap_{n\geq 1} \overline{D_{s_1s_2\dots s_n}} \  \ \ \text{ and } \ \ \  \Sigma' = \bigcap_{n\geq 1} \overline{D_{s_1's_2'\dots s_n'}}. \]
	Now, consider any lift $\tilde D_1$ of $D_{s_1}$; by Lemma \ref{lem:restriction}, $\pi\colon\tilde D_1\to D_{s_1}$ is a universal covering, but since $D_{s_1}$ is (by hypothesis) not simply connected, this covering has infinite degree. In particular, $\partial\tilde D_1$ consists of infinitely many arcs with endpoints in $\partial\DD$, one of which (say $C_1$) separating $0$ from all the others (see Lemma \ref{lemma-lifts-of-curves}). The domain $\tilde D_1$ also contains infinitely many lifts of $\sigma_{s_1}\subset\partial D_{s_1}$; we pick any one of these distinct from $C_1$ and cross back into a component $\tilde \Omega_1$ of $\pi^{-1}(\Omega_1)$. Once there, we choose a crosscut $C_1'$ corresponding to a preimage of $\sigma_{s_1'}\subset \partial  D_{s_1'}$. After crossing $C_1'$ into a component $\tilde D_1'$ of $\pi^{-1}(D_{s_1'})$, which contains (once again) infinitely many preimages of $\sigma_{s_1'}$; we can cross one that is distinct from $C_1'$, then cross another component of $\pi^{-1}(\sigma_{s_1})$ into a new component of $\pi^{-1}(D_{s_1})$. This time, we choose a preimage $C_2$ of $\sigma_{s_1s_2}$, which will take us into a preimage of $D_{s_1s_2}$. Repeating this process inductively, we produce sequences $C_n$ and $C_n'$ of fundamental crosscuts of $\DD$, which are shrinking to a point $e^{i\theta}\in\partial\DD$. This point, by construction, has both $\underline{s}$ and $\underline{s}'$ as associated boundary sequences.

	There are uncountably many points with this property because we made countably many choices among countably many possibilities at each choice.
\end{proof}

\end{document}
